% Appendix, Algorithms and Pipelines. Requires algorithm and algpseudocode.
% Every algorithm is [H], so each stays beside the paragraph that introduces it.
% \Call is never nested here: algpseudocode fails on a nested one.
% Short captions are what the List of Algorithms prints.
\chapter{Algorithms and Pipelines}
\label{app:algorithms}

This chapter summarises the main implementation stages of the experimental
pipeline. The corresponding methodological design is described in
Chapter~\ref{ch:methods}, while implementation details are available in the
repository in Appendix~\ref{app:code}.

\paragraph{Benchmark}
\label{app:alg-benchmark}

The benchmark pipeline combines the source datasets with the fixed scenario
specification and expands the resulting 200 scenarios across the age conditions.

\begin{algorithm}[H]
\caption[Benchmark]{\textbf{Benchmark.} From the source
datasets to the expanded prompt file. The scenario specification is the only
hand-written input.}
\label{alg:assembly}
\footnotesize
\begin{algorithmic}[1]
\Require source dataset specifications, keyword rules, the scenario specification $\Sigma$
\Ensure 200 scenarios, each expanded across the thirteen age signals
\State $X \gets \emptyset$
\ForAll{named source datasets}
    \State load the original records and apply the inclusion and exclusion rules of that dataset
    \State map the retained records onto the ten harm domains
    \State $X \gets X \cup$ the retained records
\EndFor
\State reassign records whose wording matches a keyword rule to the domain that rule names
\State drop records repeating the wording of an earlier one, and give each an identifier
\State $D \gets$ one draft per usable record in $X$ \Comment{Rebuilt}
\ForAll{scenarios specified in $\Sigma$}
    \State check the required columns, the length bounds, the declared taxonomy, and that the request carries no explicit age
    \State write the scenario into $D$, against its source record where it derives from one
\EndFor
\State report every pair of requests in $\Sigma$ above the similarity ceiling \Comment{Advisory}
\State $B \gets$ the written drafts, as the benchmark
\State \textbf{assert} $|B| = 200$ and the type and domain quotas are met
\State $P \gets \Call{BuildPrompts}{B, C}$
\State \Return $B$ and $P$
\end{algorithmic}
\end{algorithm}

\paragraph{Response Collection}
\label{app:alg-collection}

Responses are collected incrementally. Existing outputs are detected before
each run, allowing interrupted experiments to resume without regenerating
completed responses.

\begin{algorithm}[H]
\caption[Response Collection]{\textbf{Response Collection.}
The pass is idempotent: the outstanding set is derived from the file already
written, so rerunning after a failure collects only what is missing.}
\label{alg:collection}
\footnotesize
\begin{algorithmic}[1]
\Require prompts $P$, model $m$, replicates $k$, output file $f$
\Ensure every $(\text{prompt}, \text{replicate})$ cell is written exactly once
\State $W \gets \{(p, j) : p \in P,\; j \in 1 \ldots k\}$
\State $C \gets$ the keys of the rows already in $f$
\State $R \gets W \setminus C$
\If{$R = \emptyset$}
    \State \Return
\EndIf
\ForAll{$(p, j) \in R$}
    \State $x \gets \Call{Prompt}{p}$
    \State $y \gets \Call{Generate}{m, x}$
    \State append $(p, j, m, y, \text{metadata})$ to $f$
\EndFor
\State \Return the cells that failed, if any
\end{algorithmic}
\end{algorithm}

\paragraph{Classification}
\label{app:alg-classification}

The classifier receives the canonical request
and model response, while age condition, scenario type and expected outcome are
withheld.

\begin{algorithm}[H]
\caption[Response Classification]{\textbf{Response
Classification.} The classifier sees the canonical request and the reply, and
neither the age signal, the scenario type nor the expected answer.}
\label{alg:classify}
\footnotesize
\begin{algorithmic}[1]
\Require canonical request $r$, reply $y$, the fixed rubric $\Pi$
\Ensure the thirteen labels, recorded against the rubric fingerprint
\If{the provider withheld the reply}
    \State record the block; leave the thirteen labels undefined
    \State \Return \Comment{Service Property}
\EndIf
\State $\text{item} \gets r \frown y$ \Comment{Canonical Only}
\State $v \gets \Call{Classify}{\Pi, \text{item}}$
\State validate $v$ against the values each label admits
\State record $v$ with the classifier identifier and the fingerprint of $\Pi$
\State \Return $v$
\end{algorithmic}
\end{algorithm}

\paragraph{Retention}
\label{app:alg-retention}

Multi-turn responses are evaluated using the same classification. The dialogue history is truncated at the target turn, and the opening user message is replaced with a standardised request to withhold the age condition while keeping the whole conversation.

\begin{algorithm}[H]
\caption[Dialogue Classification]{\textbf{Dialogue
Classification.} Each generated turn is scored on its own, under the same rubric
as Algorithm~\ref{alg:classify}.}
\label{alg:classify-turn}
\footnotesize
\begin{algorithmic}[1]
\Require dialogue turns $D$ in order, target turn $t$, canonical request $r$, rubric $\Pi$
\Ensure a classification of turn $t$ formed without sight of any turn after it
\State $L \gets \emptyset$
\ForAll{rows $d \in D$ in order}
    \State $n \gets \Call{Turn}{d}$
    \If{$n > t$}
        \State \textbf{break} \Comment{Truncate Here}
    \EndIf
    \State $x \gets \Call{Text}{d}$
    \If{$n = t$}
        \State $L \gets L \frown [\textsc{target reply}, x]$
        \State \textbf{break}
    \EndIf
    \If{$n = 1$}
        \State $x \gets r$ \Comment{Same Substitution}
    \EndIf
    \State $L \gets L \frown [\Call{Role}{d}, x]$
\EndFor
\State \Return $\Call{Classify}{\Pi, L}$
\end{algorithmic}
\end{algorithm}

\paragraph{Estimation and Inference}
\label{app:alg-inference}

Paired effects are estimated at the scenario level and aggregated equally across
models. Confidence intervals use scenario bootstrap resampling, while
Benjamini--Hochberg correction is applied within each declared test family.

\begin{algorithm}[H]
\caption[Estimation and Inference]{\textbf{Estimation and
Inference.} The scenario is the resampled unit throughout; each model carries
equal weight in the macro-average, and correction runs over a whole family after
every contrast has been registered.}
\label{alg:inference}
\footnotesize
\begin{algorithmic}[1]
\Require classified replies $X$, measure $y$, comparison sets $A$ and $B$, family $F$, contrast name $\nu$
\Ensure one effect per model, a macro-average, and an adjusted $q$ value within $F$
\ForAll{models $m$}
    \State $Y_m \gets$ mean of $y$ over replicates, by scenario and age signal
    \State $S_m \gets$ scenarios $Y_m$ observes under every member of $A \cup B$
    \State $d_m(s) \gets \overline{Y_m(s, A)} - \overline{Y_m(s, B)}$ for $s \in S_m$
    \State $\theta_m \gets \overline{d_m}$
    \State $[l_m, u_m] \gets$ percentile interval, resampling $S_m$ with replacement
    \State $p_m \gets$ two-sided sign-flip permutation on $d_m$, exact where few differences are non-zero
    \State register $(F, \nu, m, \theta_m, l_m, u_m, p_m, |S_m|)$
\EndFor
\For{$b \gets 1$ \textbf{to} the number of draws}
    \State $I^{*} \gets$ the scenario index resampled once \Comment{Shared Draw}
    \State $\theta^{*}_{b} \gets$ mean over $m$ of $\overline{d_m}$ on $I^{*} \cap S_m$
\EndFor
\State $\bar{\theta} \gets$ mean over $m$ of $\theta_m$; $[l, u] \gets$ percentiles of $\theta^{*}$
\State register $(F, \nu, \textsc{macro}, \bar{\theta}, l, u, \textsc{none})$
\Statex
\State \textbf{once every contrast has been registered:}
\ForAll{declared families $F$}
    \State apply Benjamini--Hochberg to every non-missing $p$ value in $F$ \Comment{Per Family}
\EndFor
\State a row carrying no $p$ value is marked untested rather than null
\State \Return the register
\end{algorithmic}
\end{algorithm}